\chapter{HAAR MEASURE. CONVOLUTION.}

The notions of length, area and volume among the Greeks are essentially based on their invariance under displacements: ``Things which coincide (ἐφαρμόζοντα) are equal'' (Eucl. El., Book I, ``Common notions'' 4); and it is by an ingenious use of this principle that all the formulae are obtained giving the areas or volumes of the classical ``figures'' (polygons, conics, polyhedra, spheres, etc.), sometimes by a procedure of finite decomposition, sometimes by ``exhaustion''.¹ In modern language, it can be said that what was done by the Greek geometers, is to prove the existence of ``functions of sets'' which are additive and invariant under displacements, but defined only for sets of a very particular type. The integral calculus can be considered as replying to the need for enlarging the domain of definition of these functions of sets, and, from Cavalieri to H. Lebesgue, it is this preoccupation that will be in the forefront of the research of the analysts; as for the property of invariance under displacements, it becomes a secondary consideration, having become a trivial consequence of the general formula for change of variables in double or triple integrals and because of the fact that an orthogonal transformation has determinant equal to ±1. Even in the non-Euclidean geometries (where however the group of displacements is different), the point of view remains the same: in a general way, Riemann defines the infinitesimal elements of area or volume (or their analogues for dimensions ≥ 3) starting with a \(ds^2\) by means of the classical Euclidean formulae, and their invariance under transformations that leave invariant the \(ds^2\) is thus almost a tautology.

It is only around 1890 that other less immediate extensions of the notion of measure invariant under a group appear, with the development of the theory of the integral invariants, notably by H. Poincaré and E. Cartan;

H. Poincaré only foresees groups with one parameter operating on a part of space, whereas E. Cartan is mainly interested in groups of displacements, but operating in other spaces than those in which they are defined. For example, he determines thus among others ({[}52 a{]}, v. II \(_{1}\) , pp.~265-302) the invariant (under a group of displacements) measure in the space of straight lines in R \(^{2}\) or in R \(^{3}\) ; \(^{2}\) further, he points out that in general the integral invariants for a Lie group are none other than particular differential invariants, and that it is thus possible to determine all of them by the methods of Lie. It seems however that no one thought to consider or to use a measure invariant on the group itself before the fundamental work of A. Hurwitz in 1897 {[}168{]}. Seeking to construct polynomials (on R \(^{n}\) ) invariant under the orthogonal group, Hurwitz starts from the remark that, for a finite group of linear transformations, the problem is solved immediately by taking the mean of the transforms s.P of an arbitrary polynomial P by all the elements s of the group; which gives him the idea of replacing, for the orthogonal group, the mean by an integral relative to an invariant measure; he gives explicitly the expression for this latter by means of the parametric representation by the angles of Euler, but observes immediately (independently of E. Cartan) that the methods of Lie supply the existence of an invariant measure for all Lie groups. Perhaps because of the decline of the theory of invariants at the beginning of the XXth century, the ideas of Hurwitz had hardly any immediate reaction, and would not be considered valuable until after 1924, with the extension to compact Lie groups, by I. Schur and H. Weyl, of the classical theory of Frobenius (p.~121) on the linear representations of finite groups. The first limits himself to the case of the orthogonal group, and shows how the method of Hurwitz allows the extension of the classical relations of orthogonality of characters; an idea that H. Weyl combines with the work of E. Cartan on semisimple Lie algebras to obtain the explicit expression for the characters of the irreducible representations of compact Lie groups and the complete reducibility theorem {[}331 b{]}, then, by a bold extension of the notion of ``regular representation'', the famous theorem of Peter-Weyl, a perfect analogy of the decomposition of the regular representation into its irreducible components in the theory of finite groups {[}332{]}.

One year previously, O. Schreier had founded the general theory of topological groups {[}276 a{]}, and from that moment it was clear that the arguments of the memoir of Peter-Weyl would remain valid as such for every topological group on which could be defined an ``invariant measure''. Truth to tell, the general notions of topology and measure were still at this time in full formation, and neither the category of topological groups on which one could hope to define an invariant measure, nor the sets for which this ``measure'' would be defined, seemed clearly delimited. The only evident point was that one could not hope to extend to the general case the infinitesimal methods proving the existence of an invariant measure on a Lie group. But, another stream of ideas, issuing from work on the Lebesgue measure, was leading precisely to more direct methods of attack. Hausdorff had proved, in 1914, that there does not exist an additive function on sets not identically zero defined for all subsets of \(\mathbf{R}^3\) and invariant under displacements, and it was natural to find out if this result was still valid for \(\mathbf{R}\) and \(\mathbf{R}^2\) : a problem that was solved by S. Banach in 1923 in a surprising way, by showing on the contrary that in these two cases such a ``measure'' well and truly existed {[}14 b{]}; his procedure, very ingenious, relied already on a construction by transfinite induction and on the consideration of ``means'' \(\frac{1}{n} \sum_{k=1}^{n} f(x + \alpha_k)\) of translates of a function by elements of the group. \(^3\) It is analogous ideas that allowed A. Haar, in 1933 {[}139{]}, to make the decisive step by proving the existence of an invariant measure for locally compact groups with a countable base of open sets: inspired by the method for the approximation of a volume, in classical integral Calculus, by a juxtaposition of congruent cubes with arbitrarily small sides, he obtains, with the help of the diagonal procedure, the invariant measure as the limit of a sequence of ``approximating measures'', a procedure that still is essentially the one used today. This discovery had a very big repercussion, in particular because it allowed J. von Neumann to solve immediately, for compact groups, Hilbert's famous ``5th problem'' on the characterisation of Lie groups by purely topological properties (to the exclusion of all previously given differential structure). But it was noticed immediately that, in order to be able to use the invariant measure in an efficacious way, it was not enough only to know of its existence, but also to know that it was unique up to a factor; this point was first proved by J. von Neumann for compact groups, by using a method for the definition of the Haar measure by ``means'' of continuous functions, analogous to those of Banach {[}324 e{]}; then J. von Neumann {[}324 f{]} and A. Weil {[}330 c{]}, by different methods, obtained simultaneously the uniqueness in the case of locally compact groups, A. Weil indicating at the same time how the procedure of Haar could be extended to general locally compact groups. It is also A. Weil (loc. cit.) who obtained the condition for the existence of a relatively invariant measure on a homogeneous space, and showed finally that the existence of a ``measure'' (endowed with reasonable properties) on a separated topological group, implied ipso facto that the group is locally precompact. These works completed essentially the general theory of Haar measure; the only more recent addition that needs quoting is the notion of quasi-invariant measure, which hardly emerged until around 1950, linked with the theory of representations of locally compact groups in Hilbert spaces.

The history of the convolution product is more complex. Already from the beginning of the XIXth century, it is noticed that if, for example, \(F(x,t)\) is an integral of a partial differential equation in x and t, linear and with constant coefficients, then

\[
\int_ {- \infty} ^ {+ \infty} F (x - s, t) f (s) d s
\]

is also an integral of the same equation; Poisson, among others, already before 1820, uses this idea to write the integrals of the heat equation in the form

\[
\int_ {- \infty} ^ {+ \infty} \exp \left(- \frac {(x - s) ^ {2}}{4 t}\right) f (s) d s.\tag{23.1}
\]

A little later, the expression

\[
\frac {1}{2 \pi} \int_ {- \pi} ^ {+ \pi} \frac {\sin \frac {2 n + 1}{2} (x - t)}{\sin \frac {x - t}{2}} f (t) d t\tag{23.2}
\]

for the partial sum of a Fourier series and the study, made by Dirichlet, of the limit of this integral when n tends to +∞, gives the first example of ``regularisation'' \(f \mapsto \rho_{n} * f\) on the torus T (in actual fact by a sequence of non-positive ``kernels'' \(\rho_{n}\) , which complicates considerably its study); under the name of ``singular integrals'', the analogous integral expressions will be a favourite subject for the analysts of the end of the XIXth century and the beginning of the XXth century, from P. du Bois-Reymond to H. Lebesgue. On R, Weierstrass uses the integral (23.1) for a proof of his theorem on approximation by polynomials, and gives in this connection the general principle of regularisation by a sequence of positive ``kernels'' \(\rho_{n}\) of the form \(x \mapsto c_{n} \rho(nx)\) . On T, the most famous example of a regularisation by positive kernels is given a little later by Fejér, and from that moment, it is the standard procedure which will be at the basis of the majority of ``summation methods'' of series of functions.

However, this work, because of the lack of symmetry of the roles that are played by the ``kernel'' and the regularised function, hardly made apparent the algebraic properties of the convolution product. It is to Volterra especially that is owed the stressing of this point. He studies in a general way the ``composition'' \(F * G\) of two functions of two variables

\[
(F * G) (x, y) = \int_ {x} ^ {y} F (x, t) G (t, y) d t
\]

that he sees as a generalisation, by ``passing from the finite to the infinite'', of the product of two matrices {[}322 a{]}. He distinguishes very early on the case (``of the closed cycle'' so called because of its interpretation in the theory of heredity) where F and G depend only on y - x; it is then the same with \(H = F * G\) , and if we set \(F(x, y) = f(y - x)\) , \(G(x, y) = g(y - x)\) , we have

\[
H (x, y) = h (y - x),
\]

with

\[
h (t) = \int_ {0} ^ {t} f (t - s) g (s) d s
\]

so that, for \(t \geq 0\) , \(h\) coincides with the convolution of functions \(f_1, g_1\) equal respectively to \(f\) and \(g\) for \(t \geq 0\) , and to 0 for \(t < 0\) .

Meanwhile, the algebraic formalism developed by Volterra did not bring out the links with the group structure of R and the Fourier transform. We do not need to write the history of this latter here; but it is appropriate to note that starting with Cauchy the analysts who deal with the Fourier integral are mainly concerned with finding wider and wider conditions for the validity of the different formulae of ``inversion'', and neglect somewhat its algebraic properties. It is certainly not possible to say as much of the work of Fourier himself (or of that of Laplace on the analogous integral \(\int_{0}^{+\infty} e^{-st} f(t) dt\) ); but these transformations had been introduced essentially in connection with linear problems, and it is therefore not very surprising that for a long time it had not been thought worth considering the product of two Fourier transforms (with the exception of the products of trigonometric series or of entire series, but the link with the convolution of discrete measures could obviously not have been noticed in the XIXth century). The first mention of this product and of the convolution on R is probably to be found in a memoir of Tchebychef {[}306{]}, in connection with questions about the calculation of probabilities. Indeed, in this theory, the convolution \(\mu * \nu\) of two ``probability laws'' on R (positive measures of total mass 1) is none other than the ``product'' probability law of \(\mu\) and \(\nu\) (for the sum of the corresponding ``random variables''). Of course, with Tchebychef, it is still only a matter of the convolution of probability laws having a density (with respect to the Lebesgue measure), hence of the convolution of functions; besides it only occurs with him in an episodic way, and it will be thus in the few rare works where it appears before the period 1920-1930. In 1920, P.-J. Daniell, in a note rarely consulted {[}76 c{]}, defines the convolution of two arbitrary measures on R and the Fourier transform of such a measure, and observes explicitly that the Fourier transform takes the convolution to the ordinary product; a formalism that, starting in 1925, is going to be intensively used by the probabilists, following P. Lévy especially. But the fundamental importance of the convolution in the theory of groups is only fully recognised by H. Weyl in 1927; he realises that, for a compact group, the convolution of functions plays the role of multiplication in the group algebra of a finite group, and allows him as a consequence to define the ``regular representation''; at the same time, he finds in the regularisation the equivalent of the unit element of the algebra of a finite group. It remained to make the synthesis of all these points of view, which is accomplished in the book of A. Weil {[}330 c{]}, a prelude to subsequent generalisations that will constitute, on the one hand the normed algebras of I. Gelfand, and on the other the convolution of distributions.

The Haar measure and convolution have rapidly become essential tools in the tendency to algebraisation that distinguishes so strongly modern Analysis; we will have to mention numerous applications of it in subsequent notes. We will limit ourselves here to pointing out that one which concerns the ``variation'' of closed subgroups (and notably discrete subgroups) of a locally compact group. This theory, starting with a result of K. Mahler in the Geometry of numbers, was inaugurated in 1950 by C. Chabauty, and has just been considerably developed and deepened by Macbeath and Swierczkowski {[}212{]}.

\chapter{INTEGRATION IN NON LOCALLY COMPACT SPACES.}

If the study of the links between topology and measure theory goes back to the beginnings of the modern theory of functions of real variables, it is only very recently that integration in discrete topological spaces has been set up in a general way. Before giving the history of the work which preceded the actual synthesis, we will recollect some of the steps in the evolution of the ideas concerning the relations between topology and measure.

For Lebesgue, it is only a question of integrating the functions of one or several real variables. In 1913, Radon defined general measures on \(\mathbf{R}^n\) and the corresponding integrals; this theory is expounded in detail in the work {[}82{]} of Ch. de la Vallée Poussin and relies in a constant way on the topological properties of Euclidean spaces. A bit later, in 1915, Fréchet defines in {[}115 c{]} ``abstract'' measures on a set furnished with a \(\sigma\) -algebra and the integrals with respect to these measures; he notes that it is possible to establish thus the main results of Lebesgue's theory without using any topological means. He justifies his undertaking with the following words, taken from the introduction to {[}115 d{]}, first part: ``That for example in a space with an infinity of co-ordinates where various applications of Analysis had led to various non equivalent definitions of a convergent sequence, one of these definitions is replaced by another, nothing will be changed in the properties of the families and additive functions of sets in these spaces''. The research of Fréchet is completed by Carathéodory, to whom is due the important theorem on the extension of a function of sets into a measure. The beginning of the book of Saks {[}268{]} offers a condensed exposition of this point of view.

The discovery of Haar measure on locally compact groups (pp.~231 ff.) and the numerous applications that it receives immediately, then the work of Weil and Gelfand in Harmonic Analysis, lead in 1940 to a deep modification of this point of view: in this kind of question, it is most convenient to consider a measure to be a linear form on a space of continuous functions. This method forces a restriction to compact or locally compact spaces, but it is not a hindrance for almost the totality of the applications; even better, the introduction of Harmonic Analysis on the p-adic groups and the adèle groups by J. Tate and A. Weil allowed a spectacular renewing of analytic Number theory.

It is from quite another direction that the necessity comes to enlarge this point of view by the consideration of measures on topological spaces which are not locally compact: little by little, Probability Theory leads to the study of such spaces and supplies numerous non trivial examples. Perhaps the reason for the late influence of these developments on measure theory must be sought in the relative isolation of Probability Theory, which remained on the margin of traditional mathematical disciplines until recent times.

\section*{MEASURES ON SEQUENCE SPACES.}
\phantomsection
\addcontentsline{toc}{section}{MEASURES ON SEQUENCE SPACES.}

One of the most developed branches of classical Probability Theory is that of limit theorems (the law of large numbers, tendency towards the Gauss-Laplace law, \ldots); it is a case of a deepening of the notion of statistical regularity manifested by phenomena bringing into play very large populations. The correct mathematical formulation of these problems requires the introduction of measures on the space of sequences; these spaces, which constitute the most obvious generalisation of spaces of finite dimension, are the favourite subject of the researches on ``General Analysis'' undertaken around 1920 by Fréchet, Lévy, Lusin, \ldots{} It is in any case not fortuitous that Khintchine and Kolmogoroff, the creators of the new methods in Probability Theory, are both disciples of Lusin, and that Lévy very soon turned to probabilistic problems: these constituted the touch stone of the new methods.

The first implicit occurrence of a measure on a sequence space appears in the work devoted by E. Borel in 1909 to countable probabilities {[}32 b{]}. A very original idea of Borel consisted of the application of probabilistic results that he had just obtained to the proof of properties possessed by the decimal expansion of nearly every real number contained between 0 and 1. This application rests on the following fundamental remark: let us define every real number contained between 0 and 1 by the sequence of digits in its expansion in a given base \(q\) ( \(q \geq 2\) ); if one takes at random successively the different digits from a number \(x\) , independently the ones from the others and with an equal probability \(1/q\) for \(0,1,\ldots,q-1\) , the probability that \(x\) is to be found in an interval of [0,1] is equal to the length of this interval.

In 1923, Steinhaus {[}293{]} establishes these results rigorously and describes the precise mathematical model of the unbounded sequence of selections of the kind considered by Borel: let us take \(q = 2\) to simplify matters and denote by \(I\) the set with two elements \(\{0,1\}\) ; we furnish \(I\) with the measure \(\mu\) defined by \(\mu(0) = \mu(1) = \frac{1}{2}\) ; the elements of the product space \(I^{\mathbf{N}}\) are the sequences \(\varepsilon = (\varepsilon(n))_{n \in \mathbf{N}}\) of numbers equal to 0 or 1 and the map \(\phi : \varepsilon \mapsto \sum_{n \geq 0} \varepsilon(n). 2^{-n-1}\) is, up to a countable set, a bijection of \(I^{\mathbf{N}}\) on the interval \([0,1]\) ; what is more, \(\phi^{-1}\) transforms the Lebesgue measure on \([0,1]\) into the measure \(P\) on \(I^{\mathbf{N}}\) the product of the measures \(\mu\) on each of the factors. In fact Steinhaus does not have at his disposal a construction for product measures;

he uses the existence of the quasi-bijection \(\phi\) to construct the measure P on \(I^{N}\) starting from the Lebesgue measure on [0,1], then he gives an axiomatic characterisation of P. The isomorphism thus obtained allows the translation of the language of probability into that of measure and to apply the known theorems about the Lebesgue integral.

In the same work, Steinhaus considers the random series \(\sum_{n\geq 0}\sigma_n.a_n\) , where the signs \(\sigma_{n} = \pm 1\) are chosen at random independently the ones from the others and with the same probability \(\frac{1}{2}\) ; between 1928 and 1935, he studies numerous other random series. On their side, Paley, Wiener and Zygmund consider random Fourier series \(^{1}\) of the form \(\sum_{-\infty}^{\infty}a_n\exp (2\pi i(nt + \varPhi_n))\) ; the ``amplitudes'' \(a_{n}\) are fixed, and the ``phases'' \(\varPhi_{n}\) are independent random variables uniformly spread over [0, 1]. If the analytic difficulties vary enormously from problem to problem, the translation in terms of measure theory is the same in all cases and represents an extension of the case treated by Borel and Steinhaus; it is a matter of constructing a measure on \(\mathbf{R}^{\mathbf{N}}\) , the product of a family of measures all identical to one particular positive measure \(\mu\) of mass 1 on \(\mathbf{R}\) ; for example, the preceding random Fourier series correspond to the case where \(\mu\) is the Lebesgue measure on [0, 1].

To construct such product measures, two methods can be used. The first is a direct method, set up for the first time by Daniell {[}76 b{]} in 1918; it is found again in 1934 by Jessen {[}173{]} who will make a detailed study of the case where \(\mu\) is the Lebesgue measure on \([0,1]\) . The second method is the search for tricks analogous to those of Steinhaus in order to reduce to the case of the Lebesgue measure on \([0,1]\) ; this way of proceeding had the advantage of convenience as long as there is not at our disposal a complete account of the general theory of measure, for it allows the use without new proofs of the Lebesgue theorems. \(^{2}\)

\section*{THE THEORY OF BROWNIAN MOVEMENT.}
\phantomsection
\addcontentsline{toc}{section}{THE THEORY OF BROWNIAN MOVEMENT.}

This theory occupies an exceptional place in contemporary scientific development because of the constant and fruitful exchange which it witnesses between physical problems and ``pure'' mathematics. The study of brownian movement, discovered in 1829 by the botanist Brown, has been carried on intensively in the XIXth century by numerous physicists, \(^{3}\) but the first satisfactory mathematical model was invented by Einstein only in 1905. In the simple case of a particle moving along a straight line, the fundamental hypotheses of Einstein are formulated thus: if \(x(t)\) is the abscissa of the particle at time \(t\) , and if \(t_0 < t_1 < \ldots < t_{n-1} < t_n\) , the successive displacements \(x(t_i) - x(t_{i-1})\) (for \(1 \leq i \leq n\) ) are independent gaussian random variables. It is not the place to evoke here in detail the important experimental work of J. Perrin which motivated Einstein's theory: for our purposes, it is sufficient to remember only a remark of Perrin, according to which the observation of the trajectories of brownian movement suggested irresistibly to him ``the functions without derivative of the mathematicians''. This remark was the initial spark for Wiener.

Quite another stream of ideas draws its origins from the kinetic theory of gases, developed between 1870 and 1900 by Boltzmann and Gibbs. Let us consider a gas made up of N molecules of mass m at (absolute) temperature T and let us denote by \(v_{1}, \ldots, v_{N}\) the velocities of the N molecules of the gas; the kinetic energy of the system is equal to

\[
\frac {m}{2} (\mathbf {v} _ {1} ^ {2} + \dots + \mathbf {v} _ {N} ^ {2}) = 3 N k T\tag{24.1}
\]

where k is Boltzmann's constant. According to the ideas of Gibbs, the multitude of collisions between molecules does not allow the precise determination of the velocities of the molecules, and it is convenient to introduce a probability law P on the sphere S of the space with 3N dimensions defined by the equation (24.1). The ``microcanonical'' hypothesis consists in assuming that P is the measure of mass 1 invariant under rotation on the sphere S. Besides, Maxwell's law of velocities states that the law of probability of the component of the velocity of a molecule has a Gaussian measure of variance 2kT/m. E. Borel seems to have been the first to remark in 1914 that Maxwell's law is a consequence of the hypotheses of Gibbs and of properties of the sphere when the number of molecules is very large. He considers a sphere S in Euclidean space of large dimension and the measure P of mass 1 invariant under rotation on S; using the classical methods of approximation based on the Stirling formula, he shows that the projection of P on a co-ordinate axis is approximately Gaussian. These results are made precise a little later by Gâteaux and Lévy {[}201 a{]}. Given an integer \(m \geq 1\) and a number r > 0, let us denote by \(S_{m,r}\) the set of sequences of the form \((x_{1}, \ldots, x_{m}, 0, 0, \ldots)\) with \(x_{1}^{2} + \ldots + x_{m}^{2} = r^{2}\) ; let us denote also by \(\sigma_{m,r}\) the measure of mass 1 invariant under rotation on \(S_{m,r}\) . Stated in modern language, the result of Gâteaux and Lévy is the following: the sequence of measures \(\sigma_{m,1}\) tends strictly to the unit mass at the origin \((0, 0, \ldots)\) and the sequence of measures \(\sigma_{m,\sqrt{m}}\) tends strictly to a measure \(\Gamma\) of the form

\[
d \Gamma (x _ {1}, x _ {2}, \dots) = \prod_ {n = 1} ^ {\infty} d \gamma (x _ {n})
\]

(γ is the Gaussian measure with variance 1 on R).

The preceding measure \(\Gamma\) plays the role of a Gaussian measure in infinite dimensions. It appears strongly that Lévy had confusedly hoped to define in an intrinsic manner a Gaussian measure on all Hilbert spaces of infinite dimension. In fact, as was shown by Lévy and Wiener, the measure \(\Gamma\) is invariant in a certain sense \(^{4}\) under the automorphisms of \(l^{2}\) ; unfortunately, the set \(l^{2}\) of square summable sequences \((x_{1}, x_{2}, \ldots, x_{n}, \ldots)\) is of measure zero for \(\Gamma\) . It is known today that we must be content with a Gaussian premeasure on a Hilbert space of infinite dimension. \(^{5}\)

To Wiener is owed the essential progress: if there is no reasonable Gaussian measure on a Hilbert space of infinite dimension, it is possible to construct by means of the primitive operation a measure w on a space of continuous functions starting with a Gaussian premeasure. We shall explain succinctly the initial construction of w by Wiener {[}336{]}; it is directly influenced by the relation \(\Gamma = \lim_{m \to \infty} \sigma_{m, \sqrt{m}}\) of Gâteaux and Lévy. For every integer \(m \geq 1\) , let us denote by \(H_m\) the set of functions on T = ]0,1] which are constant on each of the intervals \(\left[\frac{k-1}{m}, \frac{k}{m}\right]\) (for \(k = 1, 2, \ldots, m\) ), and \(\pi_m\) the measure of mass 1 invariant under rotation on the Euclidean sphere of radius 1 in \(R^m\) . Let \(f_m\) be the isomorphism of \(H_m\) on \(R^m\) which associates with each function taking the value \(a_k\) on the interval \(\left[\frac{k-1}{m}, \frac{k}{m}\right]\) the vector \((a_1, a_2 - a_1, \ldots, a_m - a_{m-1})\) (whence the name of ``differential space'' dear to Wiener); let us denote by \(w_m\) the measure on \(H_m\) image of \(\pi_m\) under \(f_m^{-1}\) . Wiener defines the measure w being sought as the limit of the measures \(w_m\) . Precisely, let us denote by H the set of regular functions on T, with the uniform convergence topology (we have \(H_m \subset H\) for every integer \(m \geq 1\) ); for every uniformly continuous function F bounded on H, the limit \(A\{F\} = \lim_{m \to \infty} \int_{H_m} F(x) dw_m(x)\) exists; Wiener then obtains certain majorations by a subtle analysis of the fluctuations of the items in play, and, taking up again the arguments of compactness brought out by Daniell, he shows that we are in a situation for the application of Daniell's extension theorem. We conclude the existence of a measure w carried by \(C(T)\) and such that \(A\{F\} = \int_{C(T)} F(x) dw(x)\) . Wiener can then show that the measure w corresponds to the hypotheses of Einstein, \(^{6}\) and his estimates allowed him to give a precise meaning to the remark of Perrin on functions without a derivative: the set of functions satisfying a Lipschitz condition of order \(\frac{1}{2}\) , is negligible for w (on the contrary, for every a with \(0 < a < \frac{1}{2}\) , nearly every function satisfies a Lipschitz condition of order a).

We know today numerous constructions for the Wiener measure. Thus, Paley and Wiener use random Fourier series ({[}243{]}, chap.~IX): to every real sequence \(\mathbf{a} = (a_{n})_{n \geq 1}\) and every integer \(m \geq 0\) , let us define the function \(f_{m,a}\) on ]0, 1] by

\[
f _ {m, \mathbf {a}} (t) = a _ {1} t + 2 \sum_ {k = 2} ^ {2 ^ {m + 1}} \frac {1}{\pi k} a _ {k - 1} \sin \pi k t;
\]

it can be shown that, for \(\Gamma\) -almost all sequences \(\mathbf{a}\) , the sequence of functions \(f_{m,\mathbf{a}}\) tends to a continuous function \(f_{\mathbf{a}}\) and that \(w\) is the image of \(\Gamma\) by the map (defined almost everywhere) \(\mathbf{a} \mapsto f_{\mathbf{a}}\) . Later, Lévy gave in {[}201 b and c{]} another construction. Finally Kac, Donsker and Erdős show around 1950 how to replace in Wiener's initial construction the spherical measures \(\pi_m\) on \(\mathbf{R}^m\) by more general measures. Their results establish a solid link between Wiener's measure and the limit theorems of Probability Theory; they will be completed and systematised by Prokhorov {[}254{]} in work to which we will return later.

This is not the place to analyse the numerous and important probabilistic works caused by Wiener's discovery; today, the brownian movement appears only as one of the most important examples of a Markov process. We will only mention the application made by Kac of Wiener's measure to the solution of certain parabolic partial differential equations; it is a matter there of an adaptation of the ideas of Feynman in quantum field theory --- another example of this reciprocal influence between mathematics and the problems of physics.

\section*{PROJECTIVE LIMITS OF MEASURES.}
\phantomsection
\addcontentsline{toc}{section}{PROJECTIVE LIMITS OF MEASURES.}

This is the case of a theory that was developed especially in view of the needs of Probability Theory. The problems concerning a finite sequence of random variables \(X_{1}, \ldots, X_{n}\) are solved in principle when the law \(P_{X}\) of this sequence is known: it is a positive measure of mass 1 on \(\mathbf{R}^{n}\) , such that the probability of obtaining simultaneously the inequalities \(a_{1} \leq X_{1} \leq b_{1}, \ldots, a_{n} \leq X_{n} \leq b_{n}\) is equal to \(P_{X}(C)\) where \(C\) is the closed box \([a_{1}, b_{1}] \times \ldots \times [a_{n}, b_{n}]\) of \(\mathbf{R}^{n}\) .

In practice, the measure \(P_{X}\) has discrete support or else admits a density with respect to Lebesgue measure. When it is the case of an infinite sequence \((X_{n})_{n\geq1}\) of random variables, the law \(P_{n}\) of the partial sequence \((X_{1},\ldots,X_{n})\) is known in general for every integer \(n\geq1\) ; these given laws satisfy a compatibility condition that expresses the fact that the sequence \((P_{n})_{n\geq1}\) is a projective system of measures. Until around 1920, one defined in a more or less implicit manner the probabilities of events linked to the infinite sequence by a ``natural'' passing to the limit starting with probabilities in the finite case; thus it will be assumed that the probability that a game stops is the limit, for n tending to infinity, of the probability that it stops in at most n steps. Naturally, such a theory is hardly coherent, and nothing excludes the presence of ``paradoxes'', a same probability receiving two distinct estimations according to whether it is evaluated by one or other of two procedures each as ``natural'' as the other.

Steinhaus {[}293{]} seems to have been the first to feel the necessity of considering (for the case of the toss of a coin) not only the projective system \((P_{n})_{n\geq1}\) but its limit. Shortly before, in 1919, Daniell {[}76 d{]} had proved the existence of such projective limits in general, \(^{7}\) but this result seems to have remained unknown in Europe. It is found again in 1933 by Kolmogoroff in the work {[}183{]} where this author formulates the axiomatic conception of Probability Theory. The proofs of Daniell and Kolmogoroff use a compactness argument that rests on Dini's theorem.

The Daniell-Kolmogoroff theorem leaves nothing to be desired in the case of random sequences \((X_{n})_{n\geq1}\) , but the study of random functions undertaken starting in 1935 by Kolmogoroff, Feller and Doob conceals difficulties of quite another order. Let us consider for example an interval T of R, which represents the set of instants of observations of a ``stochastic process''; the set of possible trajectories is the product space \(R^{T}\) , considered as a projective limit of partial products \(R^{H}\) , where H runs through the set of finite subsets of T; one obtains in general a projective system of measures \((\mu_{H})\) . Kolmogoroff's theorem does indeed provide a measure on \(R^{T}\) , but it is defined on a family that is notably smaller than the Borel family. \(^{8}\) A variant of Kolmogoroff's construction, which provides a measure on a topological space, is due to Kakutani {[}177{]}, and has been rediscovered many times since: one considers \(\mu_{H}\) as a measure on \(\overline{R}^{H}\) carried by \(R^{H}\),\(^{9}\) the compact space \(E=\overline{R}^{T}\) is a projective limit of the finite products \(\overline{R}^{H}\) and a measure \(\mu\) on E can be defined as a projective limit of the \(\mu_{H}\) . But this procedure has a severe inconvenience; the elements of \(\overline{R}^{T}\) do not possess any regularity property allowing the continuation of the probabilistic study of the process --- or even simply the elimination of the parasitic values \(\pm\infty\) introduced by the compactification \(\overline{R}\) of R. It can be remedied by restricting the measure \(\mu\) of \(\overline{R}^{T}\) to a particular subspace (for example \(C(T)\) in the case of brownian movement); the fundamental difficulty arises from the fact that a functional space, even of the usual type, is not necessarily \(\mu\) -measurable in \(\overline{R}^{T}\) , and even the choice of the functional space can be in question. \(^{10}\)

A decisive step was accomplished in 1956 by Prokhorov in a work {[}254{]} which exercised a determining influence on the theory of stochastic processes. By putting into a convenient axiomatic form the methods used by Wiener in the article analysed above, he establishes a general existence theorem for projective limits of measures on functional spaces.

A more restricted class of projective systems was introduced by Bochner {[}24 b{]} in 1947; it is a case of projective systems made up of real vector spaces of finite dimension and surjective linear transformations. The projective limit of such a system is identified in a natural way with the algebraic dual \(E^*\) of a real vector space \(E\) supplied with a weak topology \(\sigma(E^*, E)\) ; a corresponding projective system of measures has a limit which is a measure \(\mu\) defined for a family which is notably smaller than the Borel family of \(E^*\) . Bochner characterised completely such ``premeasures'' by their Fourier transform, which is a function on \(E\) . But this result is hardly usable in the absence of a topology on \(E\) , in which case the possibility of considering \(\mu\) as a measure on the topological dual \(E'\) of \(E\) must be examined. In an independent way, R. Fortet and E. Mourier, while seeking to generalise to random variables with values in a Banach space certain classical results of Probability Theory (the law of large numbers, the central limit theorem) brought out also the fundamental role played by the Fourier transform in these questions. But substantial progress was only made in 1956 when Gelfand {[}125 c{]} suggested that the natural framework for the Fourier transform is not that of Banach spaces or Hilbert spaces, but that of nuclear Fréchet spaces. He conjectured that every continuous function of positive type on such a space is the Fourier transform of a measure on its dual, a result established soon after by Minlos {[}222{]}. Its importance comes above all from the fact that it is applicable to distribution spaces, and that the quasi-totality of functional spaces are Borel subspaces of the space of distributions (which thus constitutes a much better setting than \(\mathbf{R}^T\) ). \(^{11}\) The theory of random distributions is an area in full expansion, and we will be content with referring the reader to the work of Gelfand and Vilenkin {[}126{]}.

The results that we have just mentioned on projective limits use the existence of topologies on the base spaces. It can be asked if there exists an analogous theory in the case of ``abstract'' measures. Von Neumann proved already in 1935 the existence of product measures in all cases, but the discovery of a counterexample by Jessen and Andersen {[}290{]} ruined the hope that every projective system of measures admits a limit. Two palliatives were discovered: in 1949, C. Ionescu-Tulcea established the existence of countable projective limits, granted the existence of suitable decompositions, \(^{12}\) a very interesting result for the study of Markov processes; besides, it became clear that the topology of the spaces does not come in except by means of the set of compact subsets. It was therefore natural to seek to axiomatise this situation inside an abstract theory, by means of the notion of a compact class of subsets of a set. This work was done in 1953 by Marczewski (who establishes by this means a theorem on abstract projective limits) and Ryll-Nardzewski (who worked on the decomposition of measures). \(^{13}\)

\section*{MEASURES ON GENERAL TOPOLOGICAL SPACES AND STRICT CONVERGENCE.}
\phantomsection
\addcontentsline{toc}{section}{MEASURES ON GENERAL TOPOLOGICAL SPACES AND STRICT CONVERGENCE.}

The study of the links between topology and measure theory was especially thought of as the study of regularity properties of measures, and in particular that of ``exterior'' regularity and of ``interior'' regularity; \(^{14}\) interior regularity is equivalent to exterior regularity on a locally compact space countable at infinity. The construction that Lebesgue gives of the measure of sets on the straight line brings out these two kinds of regularity, and the property of exterior regularity of measures on a Polish space seems to have had public notoriety around 1935. But it is only in 1940, in an article whose diffusion was slowed down by the war, that A. D. Alexandroff {[}3{]} brought out the role of interior regularity and showed that this is possessed by the measures on a Polish space; this result was found again later by Prokhorov {[}254{]} and is often attributed mistakenly to this author. It only appeared very recently that this property extended to sublinian spaces; because of this, the importance of these spaces increased greatly, all the more so that it became known that their theory could be done without a metrisability hypothesis, and that the quasi-totality of functional spaces were sublinian (and even most often lusinian). \(^{15}\)

The definition of a mode of convergence (weak or strict) for measures is done in the most convenient way by putting in duality the space of measures with a space of continuous functions. Generalising an ancient result of F. Riesz, A. A. Markhoff established in 1938 a bijective correspondence between the positive functionals on \(\mathcal{C}(X)\) and the regular measures on a compact space X. In the work {[}3{]} already quoted, A. D.. Alexandroff extends these results to the case of a completely regular space: he introduces a hierarchy on the set of positive linear forms on the space \(\mathcal{C}^{b}(X)\) of bounded continuous functions on a completely regular space X, \(^{16}\) he defines strict convergence of bounded measures and proves among others the following two theorems:

\begin{enumerate}
\def\labelenumi{\alph{enumi})}
\item
  if X is Polish, the set of linear forms on \(\mathcal{C}^{b}(X)\) corresponding to measures is closed in the weak convergence of sequences;
\item
  if a sequence of bounded measures has a strict limit, ``there is no mass fleeing to infinity'' (it is a weak form of the reciprocal of the theorem of strict convergence of Prokhorov).
\end{enumerate}

From this fusion of notions and theorems, Prokhorov will know how to extract important results for the theory of stochastic processes, and to present them in a simple and striking form. In his great work of 1956 already quoted {[}254{]}, an important part is devoted to bounded positive measures on a Polish space; in generalising a construction of Lévy, he defines on the set of positive measures of mass 1 a distance which makes it into a Polish space, then he establishes an important compactness criterion for strict convergence. Independently of Prokhorov, Le Cam {[}197{]} has obtained a certain number of compactness results for the strict convergence of measures; he puts no metrisability hypotheses on the spaces that he considers, and his results reduce to previous theorems of Dieudonné in the locally compact case.

\chapter{LIE GROUPS AND LIE ALGEBRAS.}

\section{GENESIS.}

The theory, which has been called for almost a century the ``theory of Lie groups'', was essentially set up by one mathematician: Sophus Lie.

Before starting on its history, we will summarise briefly various earlier researches that prepared its development.

\subsection{Transformations (Klein-Lie, 1869-1872).}


Around 1860, the theory of permutation groups of a finite set is developing and begins to be used (Serret, Kronecker, Mathieu, Jordan). On the other hand, the theory of invariants, then in full flight, familiarises mathematicians with certain infinite sets of geometric transformations closed under composition (notably linear or projective transformations). But, before the work of 1868 of Jordan {[}174 b{]} on the ``movement groups'' (closed subgroups of the groups of displacements in Euclidean space of 3 dimensions), it does not seem as if any conscious link had been established between these two streams of ideas.

In 1869, the young Felix Klein (1849-1925), a pupil of Plücker, becomes the friend at Berlin of the Norwegian Sophus Lie (1842-1899), a few years older, being brought together by their common interest in the ``geometry of straight lines'' of Plücker and notably the theory of complexes of straight lines (p.~133). It is around this period that Lie conceives of one of his most original ideas, the introduction of the notion of invariant in Analysis and in differential geometry; one of the sources of this is his observation that the classical methods of integration ``by quadratures'' of differential equations all rely on the fact that the equation is invariant for a ``continuous'' family of transformations. It is from 1869 that is dated the first work (drawn up by Klein) where Lie uses this idea; he studies there the ``Reye complex'' (a set of straight lines intersecting the faces of a tetrahedron in 4 points having a given cross-ratio) and the curves and surfaces having as tangents straight lines of this complex ({[}202{]}, vol.~I, Abh. V, pp.~68-77): his method relies on the invariance of the Reye complex under the commutative group with 3 parameters (a maximal torus of PGL(4, C)) leaving invariant the corners of the tetrahedron. This same idea dominates the work written together by Klein and Lie when they found themselves in Paris in the spring of 1870 ({[}182{]}, v. I, pp.~416-420); they essentially determine there the connected commutative subgroups of the projective group of the plane PGL(3, C), and study the geometric properties of their orbits (under the name of V curves or surfaces); that gives them, by a uniform procedure, properties of varied curves, algebraic or transcendental, such as \(y = cx^m\) or the logarithmic spirals. Their testimony agrees in underlining the profound impression that was produced on them by the theories of Galois and of Jordan (the commentary of Jordan on Galois had appeared in Math. Annalen in 1869; in any case, Lie had heard of the theory of Galois already in 1863). Klein, who in 1871, begins to be interested in non-Euclidean geometries, sees there the beginning of his research for a classification principle for all known geometries, research that was to lead him in 1872 to the ``Erlangen programme''. On his side, Lie in a letter of 1873 to A. Mayer ({[}202{]}, vol.~5, p.~584), dates from his stay in Paris the origin of his ideas on transformations groups, and in a work of 1871 ({[}202{]}, vol.~I, Abh. XII, pp.~153-214), he uses already the term ``transformation group'' and states explicitly the problem of the determination of all the subgroups (``continuous or discontinuous'') of GL(n, C). Truth to tell, Klein and Lie must both have had difficulty in penetrating this new mathematical universe, and Klein speaks of the ``Treatise'' of Jordan, newly appeared, as a ``book sealed with seven seals'' ({[}182{]}, v. I, p.~51); he writes elsewhere in connection with ({[}182{]}, v. I, pp.~424-459): ``It is to Lie that belongs all that is concerned with the heuristic idea of a continuous group of operators, in particular all that touches on the integration of differential equations or partial differential equations. All the notions that he developed later in his theory of continuous groups were already to be found in embryo with him, but so little elaborated however, that I had to convince him of many details, for example at the beginning even the existence of the curves V, during the course of long conversations'' ({[}182{]}, v. I, p.~415).

\subsection{Infinitesimal transformations.}

The conception of an ``infinitely small'' transformation goes back to the beginnings of the infinitesimal Calculus; it is known that Descartes discovers the instantaneous centre of rotation by assuming that ``in the infinitely small'' every plane movement can be assimilated into a rotation; the elaboration of analytic Mechanics, in the XVIIIth century, is entirely founded on similar ideas. In 1851, Sylvester, seeking to form invariants of the linear group

GL(3, C) or of certain of its subgroups, gives to the parameters \(z_{j}\) occurring in these matrices ``infinitely small'' increments of the form \(\alpha_{j}dt\) , and expresses that a function \(f((z_{j}))\) is invariant by writing down the equation \(f((z_{j} + \alpha_{j}dt)) = f((z_{j}))\) ; this gives him a linear equation for f with partial derivatives Xf = 0, where

\[
X f = \sum_ {j} \alpha_ {j} \frac {\partial f}{\partial z _ {j}},\tag{25.1}
\]

X being thus a differential operator, ``a derivative in the direction of directed parameters \(\alpha_{j}\) '' ({[}304{]}, vol.~3, pp.~326 and 327); Sylvester seems to feel that there is there a general principle of fairly wide scope, but does not seem to have returned to the question. A little later, Cayley ({[}58{]}, v. II, pp.~164-178) proceeds in the same way for invariants of SL(2, C) in certain representations of this group and shows that it is the solution of two partial differential equations of the first order Xf = 0, Yf = 0, where X and Y are obtained as above starting from ``infinitely small'' transformations

\[
\left( \begin{array}{c c} 0 & 0 \\ d t & 0 \end{array} \right) \text {   and   } \left( \begin{array}{c c} 0 & d t \\ 0 & 0 \end{array} \right).
\]

In modern terms, this is explained by the fact that X and Y generate the Lie algebra \(\mathfrak{sl}(2,\mathbf{C})\) ; besides Cayley calculates explicitly the bracket XY-YX and shows that it also comes from an ``infinitely small'' transformation.

In his memoir of 1868 on groups of movements {[}174 b{]}, Jordan uses from beginning to end the concept of ``infinitely small transformation'', but exclusively from a geometric point of view. There is no doubt that it is with him that the idea of a group with one parameter ``generated'' by an infinitely small transformation appears: for Jordan, it is the set of transformations obtained by ``repeating suitably'' the infinitely small transformation (loc. cit. p.~243). Klein and Lie, in their memoir of 1871, use the same expression ``repeated infinitely small transformation'' ({[}182{]}, v. I, pp.~424-459), but the context shows that they mean by that an integral of a differential system. If the group with one parameter that they consider is made up of transformations \(x' = f(x, y, t)\) , \(y' = g(x, y, t)\) , the corresponding ``infinitely small transformation'' is given by

\[
d x = p (x, y) d t, d y = q (x, y) d t
\]

where \(p(x,y)=\frac{\partial f}{\partial t}(x,y,t_{0}),q(x,y)=\frac{\partial g}{\partial t}(x,y,t_{0})\) , where \(t_{0}\) corresponds to the identical transformation of the group. As Klein and Lie know the functions f and g explicitly, they have no problem in checking that the functions

\[
t \mapsto f (x, y, t) \text {   and   } t \mapsto g (x, y, t)
\]

give in a parametric form the integral curve of the differential equation

\[
q (\xi , \eta) d \xi = p (\xi , \eta) d \eta
\]

going through the point \((x,y)\) , but do not give any general reason for that; besides they do not use this fact again in the rest of their memoir.

\subsection{Contact transformations.}

In the following two years, Lie seems to abandon the theory of transformation groups (although he remains in very close contact with Klein, who publishes in 1872 his ``Programme'') in order to study contact transformations, the integration of partial differential equations of the first order and the relations between these two theories. We do not have to spell out the history of these questions here, and we will limit ourselves to mentioning a few points that seem to have played an important role in the genesis of the theory of transformation groups.

The notion of contact transformation generalises at the same time the point transformations and the transformations by reciprocal polars. Grosso modo, a contact transformation \(^{1}\) in \(C^{n}\) is an isomorphism of an open set \(\Omega\) of the variety \(T^{\prime}(C^{n})\) of cotangent vectors to \(C^{n}\) onto another open set \(\Omega^{\prime}\) of \(T^{\prime}(C^{n})\) transforming the canonical 1-form of \(\Omega\) into that of \(\Omega^{\prime}\) . In other words, if \((x_{1},\ldots,x_{n},p_{1},\ldots,p_{n})\) describes the canonical co-ordinates of \(T^{\prime}(C^{n})\) , a contact transformation is an isomorphism \((x_{i},p_{i})\mapsto(X_{i},P_{i})\) satisfying the relation \(\sum_{i=1}^{n}P_{i}dX_{i}=\sum_{i=1}^{n}p_{i}dx_{i}\) . Such transformations occur in the study of partial differential equations of the form

\[
F \left(x _ {1}, x _ {2}, \dots , x _ {n}, \frac {\partial z}{\partial x _ {1}}, \dots , \frac {\partial z}{\partial x _ {n}}\right) = 0\tag{25.2}
\]

Lie becomes familiar during the course of his research on these questions with the handling of Poisson parentheses

\[
(f, g) = \sum_ {i = 1} ^ {n} \left(\frac {\partial f}{\partial x _ {i}} \frac {\partial g}{\partial p _ {i}} - \frac {\partial g}{\partial x _ {i}} \frac {\partial f}{\partial p _ {i}}\right)\tag{25.3}
\]

and brackets \(^{2}\) \([X,Y]=XY-YX\) of differential operators of type (25.1); he interprets the Poisson parenthesis (25.3) as the effect on f of a transformation of type (25.1) associated with g, and observes on this occasion that the Jacobi identity for Poisson parentheses means that the bracket of differential operators corresponding to g and h is associated with the parenthesis \((g,h)\) . The search for functions g such that \((F,g)=0\) , which occurs in the method of Jacobi for integrating the partial differential equation (25.2), becomes for Lie that for an infinitesimal contact transformation leaving the given equation invariant. Finally, Lie is led to study the set of functions \((u_{j})_{1\leq j\leq m}\) of the \(x_{i}\) and the \(p_{i}\) such that the parentheses \((u_{j},u_{k})\) are functions of the \(u_{h}\) , and calls these sets ``groups'' (already considered in substance by Jacobi).

\section{CONTINUOUS GROUPS AND INFINITESIMAL TRANSFORMATIONS.}

Suddenly, in the autumn of 1873, Lie takes up again the study of transformation groups and obtains decisive results. For as much as it is possible to follow the unravelling of his thoughts in a few letters to A. Mayer from the years 1873-1874 ({[}202{]}, vol.~5, pp.~585-608), he starts from a ``continuous group'' of transformations on n variables

\[
x _ {i} ^ {\prime} = f _ {i} (x _ {1}, \dots , x _ {n}, a _ {1}, \dots , a _ {r}) (1 \leq i \leq n)\tag{25.4}
\]

depending effectively \(^{3}\) on r parameters \(a_{1},\ldots,a_{r}\) ; he observes that if the transformation (25.4) is the identity for the values \(a_{1}^{0},\ldots,a_{r}^{0}\) of the parameters, \(^{4}\) then the Taylor expansions of the \(x_{i}\) , limited to the first order:

\[
f _ {i} \left(x _ {1}, \dots , x _ {n}, a _ {1} ^ {0} + z _ {1}, \dots , a _ {r} ^ {0} + z _ {r}\right) = x _ {i} + \sum_ {k = 1} ^ {r} z _ {k} X _ {k i} \left(x _ {1}, \dots , x _ {n}\right) + \dots (1 \leq i \leq n)\tag{25.5}
\]

give a ``generic'' infinitely small transformation depending linearly on the r parameters \(z_{j}\)

\[
d x _ {i} = \left(\sum_ {k = 1} ^ {r} z _ {k} X _ {k i} (x _ {1}, \dots , x _ {n})\right) d t (1 \leq i \leq n).\tag{25.6}
\]

Proceeding as in his memoir with Klein, Lie integrates the differential system

\[
\frac {d \xi_ {1}}{\sum_ {k} z _ {k} X _ {k 1} (\xi_ {1} , \dots , \xi_ {n})} = \dots = \frac {d \xi_ {n}}{\sum_ {k} z _ {k} X _ {k n} (\xi_ {1} , \dots , \xi_ {n})} = d t,\tag{25.7}
\]

which gives him, for each point \((z_{1},\ldots,z_{r})\) , a group with one parameter

\[
t \mapsto x _ {i} ^ {\prime} = g _ {i} (x _ {1}, \dots , x _ {n}, z _ {1}, \dots , z _ {r}, t) (1 \leq i \leq n)\tag{25.8}
\]

such that \(g_{i}(x_{1},\ldots,x_{n},z_{1},\ldots,z_{r},0)=x_{i}\) for every i. He shows in an ingenious way, using the fact that the transformations (25.4) make up a set closed under composition, that the group with one parameter (25.8) is a subgroup of the given group ({[}202{]}, vol.~5, Abh. VII, pp.~32-63). The new idea, key to the whole theory, is to proceed on to the second order in the Taylor expansions of the functions (25.4). The thread of his argument is fairly confused and heuristic ({[}202{]}, vol.~5, Abh. VII, pp.~32-63 and {[}202{]}, vol.~5, pp.~600-601); it can be set out as follows. For the \(z_{j}\) that are fairly small, one can put t=1 in (25.8), and thus are obtained new parameters \(z_{1},\ldots,z_{r}\) for the transformations of the group (it is in fact the first appearance of the ``canonical parameters''). By definition, we have with regard to (25.7)

\[
\frac {\partial g _ {i}}{\partial t} = \sum_ {k} z _ {k} X _ {k i} (x _ {1} ^ {\prime}, \dots , x _ {n} ^ {\prime}),
\]

from where

\[
\frac {\partial^ {2} g _ {i}}{\partial t ^ {2}} = \sum_ {k, j} z _ {k} \frac {\partial X _ {k i}}{\partial x _ {j}} (x _ {1} ^ {\prime}, \dots , x _ {n} ^ {\prime}) \frac {\partial x _ {j} ^ {\prime}}{\partial t} =
\]

\[
\sum_ {k, j} z _ {k} \frac {\partial X _ {k i}}{\partial x _ {j}} (x _ {1} ^ {\prime}, \dots , x _ {n} ^ {\prime}) \left(\sum_ {h} z _ {h} X _ {h j} (x _ {1} ^ {\prime}, \dots , x _ {n} ^ {\prime})\right)
\]

which gives

\[
\begin{array}{l} x _ {i} ^ {\prime} = x _ {i} + \left(\sum_ {k} z _ {k} X _ {k i} (x _ {1}, \dots , x _ {n})\right) t \\ \quad + \frac {1}{2} \left(\sum_ {k, h, j} z _ {k} z _ {h} \frac {\partial X _ {k i}}{\partial x _ {j}} (x _ {1}, \dots , x _ {n}) X _ {h j} (x _ {1}, \dots , x _ {n})\right) t ^ {2} + \dots , \end{array}
\]

from where, for \(t = 1\) , the Taylor expansions with respect to the parameters \(z_{j}\)

\[
x _ {i} ^ {\prime} = x _ {i} + \left(\sum_ {k} z _ {k} X _ {k i}\right) + \frac {1}{2} \left(\sum_ {k, h, j} z _ {k} z _ {h} X _ {h j} \frac {\partial X _ {k i}}{\partial x _ {j}}\right) + \dots (1 \leq i \leq n).\tag{25.9}
\]

Let us write these relations in brief as \(x' = G(x, z)\) between vectors

\[
\boldsymbol {x} = (x _ {1}, \dots , x _ {n}), \boldsymbol {x} ^ {\prime} = (x _ {1} ^ {\prime}, \dots , x _ {n} ^ {\prime}), z = (z _ {1}, \dots , z _ {r});
\]

the fundamental property of closure of the sets of these transformations under composition is written as

\[
G (G (\boldsymbol {x}, \boldsymbol {u}), \boldsymbol {v}) = G (\boldsymbol {x}, H (\boldsymbol {u}, \boldsymbol {v}))\tag{25.10}
\]

where \(H = (H_{1},\ldots ,H_{r})\) is independent of \(x\) ; it is immediate that \(H(u,0) = u,H(0,v) = v\) , whence the expansions

\[
H _ {i} (u, v) = u _ {i} + v _ {i} + \frac {1}{2} \sum_ {h, k} c _ {i k h} u _ {h} v _ {k} + \dots ,\tag{25.11}
\]

the terms not written down being non linear either in u or in v. Transforming (25.10) with the help of (25.9) and (25.11), then comparing the terms in \(u_{h}v_{k}\) of the two members, Lie obtains the relations

\[
\sum_ {j = 1} ^ {n} \left(X _ {h j} \frac {\partial X _ {k i}}{\partial x _ {j}} - X _ {k j} \frac {\partial X _ {h i}}{\partial x _ {j}}\right) = \sum_ {l = 1} ^ {r} c _ {l h k} X _ {l i} (1 \leq h, k \leq r, 1 \leq i \leq n).\tag{25.12}
\]

His expertise with the theory of partial differential equations leads him to write these conditions in a simpler form: following the model of (25.1), he associates with each of these r infinitely small transformations obtained by putting \(z_{k}=1\) , \(z_{h}=0\) for \(h\neq k\) in (25.6), the differential operator

\[
A _ {k} (f) = \sum_ {i = 1} ^ {n} X _ {k i} \frac {\partial f}{\partial x _ {i}},\tag{25.13}
\]

and rewrites the conditions (25.12) in the form

\[
[ A _ {h}, A _ {k} ] = \sum_ {l} c _ {l h k} A _ {l},\tag{25.14}
\]

a keystone of his theory. Until then, he had used indiscriminately the terms ``infinitely small transformation'' and ``infinitesimal transformation'' (e. g. {[}202{]}, vol.~5, Abh. I, pp.~1-4); the simplicity of the relations (25.14) leads him to call the operator (25.13) the ``symbol'' of the infinitesimal transformation \(dx_{i} = X_{ki}dt\) ( \(1 \leq i \leq n\) ) ({[}202{]}, vol.~5, Abh. III, pp.~42-75) and very rapidly, it is the operator (25.13) itself that he will call ``infinitesimal transformation'' ({[}202{]}, vol.~5, Abh. III, pp.~42-75 and {[}202{]}, vol.~5, p.~589).

He then becomes aware of the tight links which unite the theory of ``continuous groups'' with his previous research on contact transformations and partial differential equations. This link fills him with enthusiasm: ``My old work was so to speak all ready in advance to found the new theory of groups of transformations'' he wrote to Mayer in 1874 ({[}202{]}, vol.~5, p.~586).

In the following years, Lie pursued the study of transformation groups. Apart from the general theorems summarised here (§ III), he obtains a certain number of more particular results: the determination of transformation groups of the straight line and of the plane, subgroups of small codimension in projective groups, groups with at most 6 parameters, etc.. For all that he does not abandon differential equations. In fact, it even appears that for him, the theory of transformation groups was to be an instrument for integrating differential equations, where the transformation group would play a role analogous to that of the Galois group of an algebraic equation. \(^{5}\) Let us note that this research leads him likewise to introducing certain sets of transformations with an infinity of parameters, which he calls ``infinite and continuous groups''; \(^{6}\) he keeps the name ``finite and continuous groups'' for the transformation groups with a finite number of parameters of type (25.4) above.

\section{	exorpdfstring{THE ``DICTIONARY'' LIE GROUPS --- LIE ALGEBRAS.}{THE “DICTIONARY” LIE GROUPS — LIE ALGEBRAS.}}

The theory of ``finite and continuous'' groups, developed by Lie in numerous memoirs starting in 1874, is set out systematically in the imposing treatise ``Theorie der Transformationsgruppen'' ({[}203{]}, 1888-1893), written in collaboration with F. Engel; \(^{7}\) it is the object of the first volume and of the last five chapters of the third, the second being consecrated to contact transformations.

As the title indicates, there was never any question in this work of anything other than transformation groups, in the meaning of the equations (25.4), where the space of ``variables'' \(x_{i}\) and the space of ``parameters'' \(a_{j}\) initially both play equally important roles. Besides the concept of ``abstract'' group is not clearly expressed at this time; when in 1883 ({[}202{]}, vol.~5, Abh. XII, pp.~311-313) Lie remarks that with the notation of (25.10), the equation \(w = H(u, v)\) , which gives the parameters of the product of two transformations of the group, defines a new group, it is as a transformation group on the space of parameters that he is considering it, obtaining thus what he calls the ``group of parameters'' (he even obtains two of them, which are none other than the group of left translations and the group of right translations). \(^{8}\)

The variables \(x_{i}\) and the parameters \(a_{j}\) in the equations (25.4) are assumed complex in principle (except in chaps. XIX-XXIV of v. 3), and the functions \(f_{i}\) analytic; Lie and Engel are of course conscious of the fact that these functions are not in general defined for all complex values of the \(x_{i}\) and the \(a_{j}\) and that, in consequence, the composition of such transformations raises serious difficulties ({[}203{]}, v. I, pp.~15-17, pp.~33-40 and passim); and although, later, they express themselves almost always as if composition of the transformations that they study were always possible without restriction, it is no doubt only for the convenience of the statements, and they re-establish explicitly the ``local'' point of view each time that it is necessary (cf.~loc. cit., pp.~168 or 189 for example or ibid., v. 3, p.~2, note at the bottom of the page); in other words, the mathematical object that they are studying is near to what we would now call a partial operation law. They do not fail to consider, on occasion, global groups, for example the 4 series of classical groups ({[}203{]}, v. 3, p.~682), but do not seem to have set themselves the question of what in general a ``global group'' can be; it suffices for them to be able to obtain, for the ``parameters'' of the classical groups (the ``variables'' of these groups do not introduce any difficulty, since it is a case of linear transformations of \(\mathbf{C}^n\) ), systems of ``local'' parameters in the neighbourhood of the identical transformation, without their being concerned about the domain of validity of the formulae that they write down. They set themselves however one problem that stands out clearly from the local theory:⁹ the study of ``mixed'' groups, that is to say groups having a finite number of connected components, such as the orthogonal group ({[}203{]}, v. I, p.~7). They present this study as that of a set of transformations closed under composition and taking inverses, which is the union of sets \(H_{j}\) such that each is described by systems of functions \(f_{i}^{(j)}\) as in (25.4); the number of (essential) parameters of each \(H_{j}\) is even assumed a priori to depend on j, but they show in fact that this number is the same for all the \(H_{j}\) . Their main result is then the existence of a finite continuous group G such that \(H_{j} = G \circ h_{j}\) for an \(h_{j} \in H_{j}\) and for all j; they establish also that G is normal in the mixed group and remark that the determination of the invariants of this latter reduces to that of the invariants of G and of a discontinuous group ({[}203{]}, v. I, chap.~18).

The general theory developed in {[}203{]} finishes up (without that being stated in a very systematic way by the authors) by forging a ``dictionary'' making the translation from properties of ``finite continuous'' groups to those of the set of their infinitesimal transformations. It is based on the ``three theorems of Lie'', each one of which is made up of an assertion and its converse.

The first theorem ({[}203{]}, v. I, pp.~33 and 72 and v. 3, p.~563) affirms in the first instance that if in (25.4) the parameters are effective, the functions \(f_{i}\) satisfy a system of partial differential equations of the form

\[
\frac {\partial f _ {i}}{\partial a _ {j}} = \sum_ {k = 1} ^ {r} \xi_ {k i} (f (x, a)) \psi_ {k j} (a) (1 \leq i \leq n)\tag{25.15}
\]

where the matrix \((\xi_{kl})\) is of maximum rank and \(\det(\psi_{kj}) \neq 0\) ; reciprocally, if the functions \(f_{i}\) have this property, the formulae (25.4) define a groupuscule of transformations.

The second theorem ({[}203{]}, v. I, pp.~149 and 158, and v. 3, p.~590) gives relations between the \(\xi_{ki}\) on the one hand, the \(\psi_{ij}\) on the other: the conditions on the \(\xi_{ki}\) can be written in the form

\[
\sum_ {k = 1} ^ {n} \left(\xi_ {i k} \frac {\partial \xi_ {j l}}{\partial x _ {k}} - \xi_ {j k} \frac {\partial \xi_ {i l}}{\partial x _ {k}}\right) = \sum_ {k = 1} ^ {r} c _ {i j} ^ {k} \xi_ {k l} (1 \leq i, j \leq r, 1 \leq l \leq n)\tag{25.16}
\]

where the \(c_{ij}^{k}\) are constants \((1 \leq i, j, k \leq r)\) anti-symmetric in i, j. The conditions on the \(\psi_{ij}\) in the form given by Maurer, are:

\[
\frac {\partial \psi_ {k l}}{\partial a _ {m}} - \frac {\partial \psi_ {k m}}{\partial a _ {l}} = \frac {1}{2} \sum_ {1 \leq i, j \leq r} c _ {i j} ^ {k} (\psi_ {i l} \psi_ {j m} - \psi_ {j l} \psi_ {i m}) (1 \leq k, l, m \leq r).\tag{25.17}
\]

By introducing the matrix \((\alpha_{ij})\) contragredient to \((\psi_{ij})\) and the infinitesimal transformations

\[
X _ {k} = \sum_ {i = 1} ^ {n} \xi_ {k i} \frac {\partial}{\partial x _ {i}}, A _ {k} = \sum_ {j = 1} ^ {r} \alpha_ {k j} \frac {\partial}{\partial a _ {j}} (1 \leq k \leq r),\tag{25.18}
\]

(25.16) and (25.17) can be written respectively:

\[
[ X _ {i}, X _ {j} ] = \sum_ {k = 1} ^ {r} c _ {i j} ^ {k} X _ {k} (1 \leq i, j \leq r)\tag{25.19}
\]

\[
\left[ A _ {i}, A _ {j} \right] = \sum_ {k = 1} ^ {r} c _ {i j} ^ {k} A _ {k}.\tag{25.20}
\]

Conversely, if given r infinitesimal transformations \(X_{k}(1 \leq k \leq r)\) , linearly independent and satisfying conditions (25.19), the one parameter subgroups generated by these transformations generate a group of transformations with r essential parameters.

Finally, the third theorem ({[}203{]}, v. I, pp.~170 and 297 and v. 3, p.~597) reduces the determination of systems of infinitesimal transformations \((X_{k})_{1\leq k\leq r}\) satisfying (25.19) to a purely algebraic problem: we must have

\[
c _ {i j} ^ {k} + c _ {j i} ^ {k} = 0\tag{25.21}
\]

\[
\sum_ {l = 1} ^ {r} (c _ {i l} ^ {m} c _ {j k} ^ {l} + c _ {k l} ^ {m} c _ {i j} ^ {l} + c _ {j l} ^ {m} c _ {k i} ^ {l}) = 0 (1 \leq i, j, k, m \leq r).\tag{25.22}
\]

Conversely, \(^{10}\) if (25.21) and (25.22) are satisfied, there exists a system of infinitesimal transformations satisfying relations (25.19), from whence a group of transformations with r parameters (in other words, the linear combinations with constant coefficients of the \(X_{k}\) make up a Lie algebra, and conversely every Lie algebra of finite dimension can be obtained in this way).

These results are completed by the study of isomorphism questions. Two groups of transformations are said to be similar if it is possible to pass from one to the other by an invertible transformation of co-ordinates on the variables and an invertible transformation of co-ordinates on the parameters: from the beginning of his research, Lie had met this notion in a natural way in connection with the definition of the ``canonical parameters''. He shows that two groups are similar if, by means of a transformation on the ``variables'', one can bring the infinitesimal transformations of the one onto those of the other ({[}203{]}, v. I, p.~329). A necessary condition for this to be so is that the Lie algebras of the two groups be isomorphic, which Lie expresses by saying that the groups are ``gleichzusammengesetzt''; but this condition is not sufficient, and a whole chapter ({[}203{]}, v. I, chap.~19) is devoted to obtaining supplementary conditions insuring that the groups should be ``similar''. On the other hand the theory of permutation groups supplied the notion of the ``holoedric isomorphism'' of two such groups (isomorphism of the underlying ``abstract'' groups); Lie transposes this notion to groups of transformations, and shows that two such groups are ``holoedric isomorphic'' if and only if their Lie algebras are isomorphic ({[}203{]}, v. I, p.~418). In particular, every group of transformations is holoedrically isomorphic to each of its groups of parameters, and that shows that, when one wishes to study the structure of the group, the ``variables'' on which they operate do not matter much and that in fact everything reduces to the Lie algebra. \(^{11}\)

Always by analogy with the theory of permutation groups. Lie introduces the notions of subgroups, normal subgroups, ``meriedric isomorphisms'' (surjective homomorphisms), and shows that they correspond to that of subalgebras, ideals and surjective homomorphisms of Lie algebras; besides he had met very early on a particularly important example of ``meriedric isomorphism'', the adjoint representation, and had recognised its links with the centre of the group ({[}202{]}, vol.~5, Abh. III, pp.~42-75). For these results, as for the fundamental theorems, the essential tool is the theorem of Jacobi-Clebsch giving the complete integrability of a differential system (one of the forms of the theorem called ``of Frobenius''); he gives a new proof of it, as it happens, using groups with one parameter ({[}203{]}, v. I, chap.~6).

The notions of transitivity and of primitivity, so important for groups of permutations, arise also naturally for the ``finite and continuous'' groups of transformations, and the treatise of Lie-Engel makes a detailed study of them ({[}203{]}, v. I, chap.~13 and passim); the relations with the subgroups stabilising a point and the notion of homogeneous space are observed (for as much as they could be without putting oneself in a global point of view) ({[}203{]}, v. I, p.~425).

Finally, the ``dictionary'' is completed, in {[}203{]}, by the introduction of the notions of derived group and soluble group (called ``integrable group'' by Lie; this terminology, suggested by the theory of differential equations, will remain in use until the work of H. Weyl) ({[}203{]}, v. I, p.~261 and v. 3, pp.~678-679); the relation between commutators and brackets had besides been observed by Lie already in 1883 ({[}202{]}, vol.~5, p.~358).

\subsection*{Other proofs of fundamental theorems.}
\phantomsection
\addcontentsline{toc}{subsection}{Other proofs of fundamental theorems.}

In {[}278 b{]}, F. Schur shows that in canonical co-ordinates the \(\psi_{ik}\) of (15) satisfy the differential equations

\[
\frac {d}{d t} (t \psi_ {i k} (t a)) = \delta_ {i k} + \sum_ {j, l} c _ {j l} ^ {k} t a _ {l} \psi_ {j i} (t a).\tag{25.23}
\]

These integrate and give a formula equivalent to the formula

\[
\varpi (X) = \sum_ {n \geq 0} \frac {1}{(n + 1) !} (\operatorname{ad} (X)) ^ {n}\tag{25.24}
\]

for the right differential \(\varpi(X)\) of the exponential map at the point X; in particular, in canonical co-ordinates, the \(\psi_{ij}\) are extended into entire functions of the \(a_{k}\) . F. Schur deduces from this a result making precise an earlier remark of Lie: if, in the definition (25.4) of groups of transformations, it is only assumed that the \(f_{i}\) are in the class \(C^{2}\) , then the group is holoedrically isomorphic to an analytic group. \(^{12}\) Following his research on the integration of differential systems, E. Cartan ({[}52 a{]}, v. II₂, p.~371) introduces in 1904 the Pfaffian forms

\[
\omega_ {k} = \sum_ {i = 1} ^ {r} \psi_ {k i} d a _ {i} (1 \leq i \leq r)\tag{25.25}
\]

(with the notation of (25.15)), later called Maurer-Cartan forms. The conditions (25.17) of Maurer can then be written

\[
d \omega_ {k} = - \frac {1}{2} \sum_ {i, j} c _ {i j} ^ {k} \omega_ {i} \wedge \omega_ {j};
\]

E. Cartan shows that it is possible to develop the theory of finite continuous groups starting with the \(\omega_{k}\) and establishes the equivalence of this point of view and that of Lie. But, for him, the interest in this method is especially that it adapts itself to ``infinite and continuous'' groups for which he pushes the theory much further that had been done by Lie, and that it allows the setting up of his theory of the generalised ``mobile marker''.

\section{THE THEORY OF LIE ALGEBRAS.}

Once the correspondence between groups of transformations and Lie algebras is acquired, the theory is going to take a clearly more algebraic turn and will be centred on a deep study of Lie algebras. \(^{13}\)

A first and short period, from 1888 to 1894, marked by the work of Engel, of his pupil Umlauf and especially Killing and E. Cartan, ended up with a series of spectacular results on complex Lie algebras. We have seen above that the notion of soluble Lie algebra was due to Lie himself, who had proved (in the complex case) the reduction theorem for soluble linear Lie algebras to triangular form ({[}203{]}, v. I, p.~270). \(^{14}\) Killing observes {[}180{]} that there exists in a Lie algebra a biggest soluble ideal (that is called the radical today), and that the quotient of the Lie algebra by its radical has null radical; he calls the Lie algebras with null radical semisimple, and proves that they are the product of simple algebras (this last notion had already been introduced by

Lie, who had proved the simplicity of the ``classical'' Lie algebras ({[}203{]}, v. 3, p.~682)).

On the other hand, Killing introduces, in a Lie algebra, the characteristic equation \(\det(\operatorname{ad}(x)-\omega.1)=0\) , already met by Lie in studying the Lie subalgebras of dimension 2 containing a given element of a Lie algebra. We refer to other historical Notes for the analysis of the methods whereby Killing, in studying in a penetrating manner the properties of the roots of the ``generic'' characteristic equation for a semisimple algebra, ends up with the most remarkable of his results, the complete determination of the simple (complex) Lie algebras. \(^{15}\)

Killing proves that the derived algebra of a soluble algebra is ``of rank 0'' (which means that ad \(x\) is nilpotent for every element \(x\) of the algebra). Soon after, Engel proves that the algebras ``of rank 0'' are soluble (this statement is in substance what is called today Engel's theorem). In his thesis, E. Cartan introduces on the other hand what is now called the ``Killing form'', and establishes the two fundamental criteria that characterise by means of this form the soluble Lie algebras and the semisimple Lie algebras.

Killing had affirmed ({[}180{]}, IV) that the derived algebra of a Lie algebra is the sum of a semisimple algebra and of its radical, which is nilpotent, but his proof was incomplete. A little later, E. Cartan announced without proof ({[}52 a{]}, v. I₁, p.~104) that more generally every Lie algebra is the sum of its radical and of a semisimple subalgebra; the only result in this direction that is indisputably established at this time is a theorem of Engel affirming the existence, in every non-soluble Lie algebra, of a simple Lie subalgebra of dimension 3. The first published proof (for complex Lie algebras) of the statement of Cartan is due to E. E. Levi {[}200{]}; another proof (equally valid for the real case) was given by J. H. C. Whitehead in 1936 {[}334 a{]}. In 1942, A. Malcev completed this result by the uniqueness theorem on the ``Levi sections'' up to conjugation.

Already from his first work, Lie had set himself the problem of the isomorphism of every Lie algebra with a linear Lie algebra. He had thought that he had solved it affirmatively by considering the adjoint representation (and so deducing a proof of his ``third theorem'') ({[}202{]}, vol.~5, Abh. III, pp.~42-75); he recognised rapidly that his proof was only correct for Lie algebras with null centre; after him, the question remained open for a long time, and was solved affirmatively by Ado in 1935 {[}2 a{]}. On the other hand, Lie had substantially set himself the problem of determining the linear representations of minimal dimension of the simple Lie algebras, and had solved it for the classical algebras; in his Thesis, Cartan had also resolved this problem for the exceptional simple algebras; \(^{16}\) the methods that he employs in this context will be generalised by him twenty years later in order to obtain all the irreducible representations of simple real or complex Lie algebras.

The property of complete reducibility of a linear representation seems to have been met for the first time (in a geometric form) by Study. In an unpublished manuscript, but quoted in ({[}203{]}, v. 3, pp.~785-788) he proves this property for the linear representation of the Lie algebra of SL(2, C), and obtains partial results for SL(3, C) and SL(4, C). Lie and Engel conjecture on this occasion that the complete reducibility theorem is valid for SL(n, C) whatever n may be. The complete reducibility of the linear representations of semisimple Lie algebras was established by H. Weyl in 1925 \(^{17}\) by an argument of a global nature (see later). The first algebraic proof was obtained in 1935 by Casimir and van der Waerden {[}55{]}; other algebraic proofs were given afterwards by R. Brauer {[}34 b{]} and J. H. C. Whitehead {[}334 b{]}.

Finally, during his research on the exponential map (cf.~infra). H. Poincaré ({[}251{]}, v. III) considers the associative algebra of differential operators of all orders, generated by the operators of a Lie algebra; he shows, substantially that, if \((X_{i})_{1\leq i\leq n}\) is a basis of a Lie algebra, the associative algebra generated by the \(X_{i}\) has as basis certain symmetric functions of the \(X_{i}\) (sums of the non-commutative ``monomials'' obtained from a given monomial by all permutations of its factors). The essential part of his proof is algebraic in nature, and allows him to obtain the structure of the enveloping algebra. Analogous proofs have been given in 1937 by G. Birkhoff {[}22 b{]} and E. Witt {[}337 b{]}. \(^{18}\)

The majority of the work quoted above is limited to real or complex Lie algebras, which are the only ones corresponding to Lie groups in the usual sense. The study of Lie algebras over a field other than R or C is tackled by Jacobson {[}172 a{]} who shows that the greater part of the classical results remain valid over a field of characteristic zero.

\section{THE EXPONENTIAL AND THE FORMULA OF HAUSDORFF.}

The first research concerning the exponential map is due to E. Study and F. Engel; Engel {[}104 b{]} remarks that the exponential is not surjective for \(\mathbf{SL}(2, \mathbf{C})\) (for example \(\left( \begin{array}{cc} -1 & a \\ 0 & -1 \end{array} \right)\) is not an exponential if \(a \neq 0\) ), but that it is for \(\mathbf{GL}(n, \mathbf{C})\) , thus also for \(\mathbf{PGL}(n, \mathbf{C})\) (this latter property had already been noted by Study for \(n = 2\) ); thus \(\mathbf{SL}(2, \mathbf{C})\) and \(\mathbf{PGL}(2, \mathbf{C})\) give an example of two groups that are locally isomorphic, but which are nonetheless very different from a global point of view. Engel shows also that the exponential is surjective in the other classical groups, augmented by homotheties; this work is taken up again and continued by Maurer, Study and others, without bringing in substantial new results.

In 1899, H. Poincaré ({[}251{]}, v. III, pp.~169-172 and 173-212), tackles the study of the exponential map from a different point of view. His memoirs seem to have been hurriedly written, for in several places he affirms that every element of a connected group is an exponential, whereas he gives examples to the contrary elsewhere. His results refer mainly to the adjoint group: he shows that a semisimple element of such a group G can be the exponential of an infinity of elements of the Lie algebra \(L(G)\) , whereas a non-semisimple element may be not an exponential. If \(\operatorname{ad}(X)\) does not have an eigenvalue which is a non-zero multiple of \(2\pi i\) , then exp is étale in X. He proves also that, if U and V describe loops in \(L(G)\) , and if W is defined by continuity such that \(e^{U}.e^{V}=e^{W}\) , we do not return necessarily to the initial determination of W. He uses a formula for residues that reduces essentially to

\[
\Phi (\operatorname{ad} X) = \frac {1}{2 \pi i} \int \frac {\Phi (\xi) d \xi}{\xi - \operatorname{ad} X}
\]

where \(\operatorname{ad}(X)\) is a semisimple element whose non-zero eigenvalues have multiplicity 1, \(\Phi\) is an entire series with a sufficiently large radius of convergence, the integral being taken over a loop enclosing the eigenvalues of ad X; he studies also what happens when X tends towards a transformation having multiple eigenvalues.

The search for expressions for \(W\) as a function of \(U\) and \(V\) in the formula \(e^U.e^V = e^W\) had already, a little before the work of Poincaré, been the subject of two memoirs of Campbell {[}46 a and b{]}. As is written a little later by Baker ``\ldots the theory of Lie suggests in an obvious way that the product \(e^Ue^V\) is of the form \(e^W\) where \(W\) is an alternating series in \(U\) and \(V\ldots\)''. The subsequent work on this subject aimed to make this assertion precise and to give an explicit formula (or a method of construction) for \(W\) (``Hausdorff formula''). After Campbell and Poincaré, Pascal, Baker {[}13{]} and Hausdorff {[}152 a{]} return to the question; each considers that the proofs of his predecessor are not convincing; the main difficulty resides in what is meant by ``alternating'': is it a case of elements of the particular Lie algebra under consideration, or of universal ``symbolic'' expressions? Neither Campbell, nor Poincaré, nor Baker express themselves clearly on this point. The memoir of Hausdorff, on the contrary, is perfectly precise; he works at first in the algebra of formal associative (non commutative) series in a finite number of indeterminates and considers U, V, W as elements of this algebra. He proves the existence of W by an argument on differential equations analogous to those of his predecessors. The same argument is used by him to prove the convergence of the series when the indeterminates in it are replaced by elements of a Lie algebra of finite dimension. As Baker had remarked, and Poincaré independently, this result can be used to give a proof of the third theorem of Lie; he clarifies the correspondence between groups and Lie algebras, for example as far as the group of commutators is concerned.

In 1947, Dynkin {[}98 a{]} takes up the question again, and obtains the explicit coefficients of the Hausdorff formula, by considering from the beginning a normed Lie algebra (of finite dimension or not, over R, C or an ultrametric field). \(^{19}\)

\section{LINEAR REPRESENTATIONS AND GLOBAL LIE GROUPS.}

None of the work of which we have been speaking tackled frankly the problem of the definition and the study of global Lie groups. The first steps along this path go back to H. Weyl. He is inspired by two theories, which until then had developed independently: that of linear representations of semisimple complex Lie algebras, due to E. Cartan, and that of linear representations of finite groups, due to Frobenius and which had just been transposed to the orthogonal group by I. Schur using an idea of Hurwitz. This latter had shown {[}168{]} how invariants can be constructed for the orthogonal groups or the unitary groups by replacing the operation of the mean on a finite group by integration relative to an invariant measure. He had also remarked that in applying this method to the unitary group, invariants for the general linear group are obtained, a first example of the ``unitarian trick''. In 1924, I. Schur {[}279 e{]} used this procedure to show the complete reducibility of the representations of the orthogonal group O(n) and of the unitary group U(n), by the construction of a positive non-degenerate invariant hermitian form; he deduces from this, by the ``unitarian trick'', the complete reducibility of the holomorphic representations of O(n, C) and of SL(n, C), establishes orthogonality relations for the characters of O(n) and of U(n) and determines the characters of O(n). H. Weyl also extends this method to complex semisimple Lie algebras {[}331 b{]}. Given such an algebra g, he shows that it possesses a ``real compact form'' (which amounts to saying that it comes from an algebra \(g_{0}\) over R, whose adjoint group \(G_{0}\) is compact, by extending the scalars from R to C). Furthermore he shows that the fundamental group of \(G_{0}\) is finite, thus that the universal cover \(^{20}\) of \(G_{0}\) is compact. He deduces from this, by an appropriate adaptation of the procedure of Schur, the complete reducibility of the representations of g, and gives also, by global means, the determination of the characters of the representations of g. In a letter to I. Schur {[}331 a{]}, H. Weyl summarises the results of Cartan, that Schur did not know (cf.~{[}279 e{]}, p.~299, note at the foot of the page) and compares the two points of view: the method of Cartan supplies all the holomorphic representations of the simply connected group of the Lie algebra g; in the case of the orthogonal group, the representations of a cover in two sheets (later called the spinor group), which had escaped Schur's notice, are also obtained; from another side, Schur's method has the advantage of proving complete reducibility and of giving the characters explicitly.

After the work of H. Weyl, E. Cartan adopts an openly global point of view in his research on symmetric spaces and Lie groups. It is this point of view that is at the basis of his exposé of 1930 ({[}52 a{]}, v. I₂, pp.~1165-1225) of the theory of ``finite and continuous'' groups. There is found in particular the first proof of the global variant of the third fundamental theorem (existence of a Lie group with given Lie algebra); Cartan shows also that every closed subgroup of a real Lie group is a Lie group, which generalises a result of J. von Neumann on the closed subgroups of the linear group {[}324 b{]}. In this Memoir, von Neumann showed also that every continuous representation of a complex semisimple group is real analytic.

After this work, the theory of Lie groups in the ``classical'' sense (that is to say of finite dimension over R or C) is more or less fixed in its main lines. The first detailed exposé of it is given by Pontrjagin in his book on topological groups {[}253{]}; he keeps there to a point of view still fairly close to that of Lie, but in distinguishing carefully the local from the global. It is followed by the book of Chevalley {[}62 d{]} which contains also the first systematic discussion of the theory of analytic varieties and of the exterior differential calculus; the ``infinitesimal transformations'' of Lie appear there as fields of vectors and the Lie algebra of a Lie group G is identified with the space of fields of vectors invariant on the left under G. He leaves to one side the ``groupuscule'' aspect and the ``groups of transformations'' aspect.

\section{EXTENSIONS OF THE NOTION OF A LIE GROUP.}

In our day, the vitality of Lie theory is manifested by the diversity of its applications (in topology, differential geometry, arithmetic, etc.) as well as by the creation of parallel theories where the structure of the underlying differential variety is replaced by a neighbouring structure (p-adic, algebraic variety, schema, formal schema, \ldots). We do not need to provide here the history of all these developments, and we will limit ourselves to two of them: Banach Lie groups and p-adic Lie groups.

\subsection{Banach Lie groups.}

It is a case of Lie groups ``of infinite dimension''. From the local point of view, a neighbourhood of 0 in an Euclidean space is replaced by a neighbourhood of 0 in a Banach space. It is what is done by G. Birkhoff in 1936 {[}22 a{]}, ending up thus with the notion of a complete normed Lie algebra and with its correspondence with a ``groupuscule'' defined on an open set of a Banach space. Around 1950, Dynkin completes these results by an extension to this case of the Hausdorff formula (cf.~supra).

The definitions and results of Birkhoff and Dynkin are local. Until recent days, it does not seem that an attempt has been made to make explicit the corresponding global theory, no doubt because of the lack of applications.
